\title{Lecture 4: RL Setting}
\author{Qi Zhang}
\date{Last Updated: September 2026}

\section{The RL Setting}
In the previous lectures, we focused on two problem settings:
the planning setting where the (finite-horizon) MDP of interest $M=(\mathcal{S},\mathcal{A}, P, R, H)$ is fully known and we aim to compute an optimal policy for it,
and the generative model setting where $(P,R)$ is unknown but can be queried for any state-action pair multiple times, and we aim to find an (near-)optimal policy after querying as few times as possible.

Obviously, the generative model setting relaxes the requirements for planing. 
In this note, we introduce the reinforcement learning (RL) setting that is even more relaxed:
$(P,R)$ is unknown and cannot be queried as a generative model; instead, we can only access them through the agent-environment interactions, i.e., {\em actually} taking actions and observing the corresponding next states and rewards.
We will focus on the infinite-horizon, discounted-reward setting where the MDP is specified by a tuple $M=(\mathcal{S},\mathcal{A},P,R,\gamma)$ with initial state distribution $d_0\in\Delta(\mathcal{S})$, and the interaction is described below, where the role of the learner's RL algorithm is highlighted in {\color{OliveGreen} green}:
\begin{protocol}
\caption{RL interaction (infinite-horizon)} \label{protocol:RL interaction}
\begin{algorithmic}[1]
\State {\color{OliveGreen} learner initializes and will maintain its internal state;}
\Repeat {~over episodes:}
    \State learner observes initial state $s_1\sim d_0$ of the current episode; 
    \For{$h=1,\ldots,H$ until termination or truncation} \Comment{{\color{Gray} truncation: $s_{H+1}$ is not terminal}}
        \State learner chooses action $a_h$ {\color{OliveGreen} based on its internal state};
        \label{protocol:RL interaction:line:choose action}
        \State learner takes $a_h$ and  observes next state $s_{h+1}\sim P(s_h,a_h)$ and reward $r_h=R(s_h,a_h)$;
        \State {\color{OliveGreen} learner updates its internal state;}
        \label{protocol:RL interaction:line:update internal state}
    \EndFor
\Until{some stopping criterion is met}
\State learner outputs a policy $\widehat{\pi}$ {\color{OliveGreen} based on its internal state};
\label{protocol:RL interaction:line:final policy}
\end{algorithmic}
\end{protocol}

In Protocol \ref{protocol:RL interaction}, the agent is referred to as the learner to emphasize the fact that we are in the RL setting.
The interaction is repeated over episodes, each of which ends either upon {\em termination} (a terminal state is reached) or upon {\em truncation} (the horizon $H$ is reached without a terminal state), and a similar interaction protocol exists for the finite-horizon case.

\paragraph{RL algorithm.}
A so-called {\em RL algorithm} essentially chooses the action to take at each timestep (line \ref{protocol:RL interaction:line:choose action}) and, to this end, maintains and updates its {\em internal state} (line \ref{protocol:RL interaction:line:update internal state});
all other lines are simply the learner observing information revealed by the environment.
In general, the learner adaptively chooses the actions based on all previous information, i.e., the decision of $a_h$ in the $k$-th episode is based on
\begin{align*}
    \underbrace{\left\{(s^j_1,a^j_1,r^j_1,\ldots,s^j_{H_j},a^j_{H_j},r^j_{H_j},s^j_{H_j+1})\right\}_{j=1}^{k-1}}_{\text{the previous $k-1$ episodes}}
    \quad \text{and } \underbrace{(s^k_1,a^k_1,r^k_1,\ldots,s^k_{h-1},a^k_{h-1},r^k_{h-1}, s^k_h)}_{\text{the in-episode transitions up to the current state}}
\end{align*}
where the superscript indexes the episodes and $H_j\leq H$ is the length of episode $j$.
The learner's internal state is thus a (possibly compressed) summary of this information.

\paragraph{RL evaluation.}
Protocol \ref{protocol:RL interaction} induces two alternative evaluation criteria for RL algorithms, both of which have been extensively studied:
\begin{itemize}
    \item {\em Exploration-exploitation.}
    The first measures the total rewards gathered {\em within} the interaction, i.e., over all episodes before the stopping criterion is met.
    Under this criterion, the learner is asked to balance between doing state-action pairs that are less tried (exploration) and redoing the state-action pairs that have been most promising so far (exploitation).

    \item {\em Pure exploration.}
    The second ignores the rewards gathered in these episodes; instead, it further asks the learner to recommend a policy $\widehat{\pi}$ by the end of the interaction (line \ref{protocol:RL interaction:line:final policy}) and measures the quality of that policy.
    Under this criterion, the learner is asked to gather high quality information within the interaction budget in order to recommend a best possible policy. 
\end{itemize}
For both criteria, the performance of an RL algorithm hinges on how efficiently it explores the environment through its action selection.

\subsection{Taxonomy of RL Algorithms}
In this note, we begin studying RL algorithms.
Crucially, in Protocol \ref{protocol:RL interaction}, the learner maintains its {\em internal state} that 
1) gets updated based on the newest transition (line \ref{protocol:RL interaction:line:update internal state})
and 2) determines the very next action to take (line \ref{protocol:RL interaction:line:choose action}).
This way, the learner is fully adaptive: in general, it chooses the current action based on all previous transitions.
The learner's internal state can include some or all of the following components:
\begin{itemize}
    \item Value estimates: $V:\mathcal{S}\to\R$, $Q:\mathcal{S}\times\mathcal{A}\to\R$
    \item Policy/actor:  $\pi:\mathcal{S}\to\Delta(\mathcal{A})$
    \item Replay buffer: $\mathcal{D}={(s,a,r,s')}$ storing previous transitions
    \item Model estimate: $(\widehat{P}, \widehat{R})$ that estimates the ground-truth $(P,R)$ based on previous transitions
\end{itemize}
RL algorithms are often characterized by what components above are (not) included in the learner's internal state.
Table \ref{table:A Taxonomy of RL Algorithms} provides a (simplified) taxonomy, although it is usually of little importance to memorize such a table. 
\begin{table}[tb]
\centering
\caption{A Taxonomy of RL Algorithms.}
\label{table:A Taxonomy of RL Algorithms}
\begin{tabular}{|l|lll|}
\hline
                       & \multicolumn{1}{l|}{\begin{tabular}[c]{@{}l@{}}Value estimates\\ $V,Q$\end{tabular}} & \multicolumn{1}{l|}{\begin{tabular}[c]{@{}l@{}}Policy\\ $\pi$\end{tabular}} & \begin{tabular}[c]{@{}l@{}}Model estimate\\ $(\widehat{P}, \widehat{R})$\end{tabular} \\ \hline
Model-free             & \multicolumn{1}{l|}{-}                                                               & \multicolumn{1}{l|}{-}                                                      & Y                                                                                     \\ \hline
Model-based            & \multicolumn{1}{l|}{-}                                                               & \multicolumn{1}{l|}{-}                                                      & Y                                                                                     \\ \hline
Value-based            & \multicolumn{1}{l|}{Y}                                                               & \multicolumn{1}{l|}{N}                                                      & N                                                                                     \\ \hline
Policy-based           & \multicolumn{1}{l|}{N}                                                               & \multicolumn{1}{l|}{Y}                                                      & N                                                                                     \\ \hline
Actor-critic           & \multicolumn{1}{l|}{Y}                                                               & \multicolumn{1}{l|}{Y}                                                      & N                                                                                     \\ \hline
Tabular                & \multicolumn{3}{l|}{Represented by tables, one entry for each $(s)$ or $(s,a)$}                                                                                                                                                                                                                                   \\ \hline
Function Approximation & \multicolumn{3}{l|}{Represented by generic functions}                                                                                                                                                                                                                                   \\ \hline
\end{tabular}
\end{table}

\section{Tabular Q-Learning}
We now introduce our first RL algorithm, {\em Q-Learning}.
Its key idea is that, in order to find an optimal policy (by the end of Protocol \ref{protocol:RL interaction} at line \ref{protocol:RL interaction:line:final policy}), it suffices to find out the optimal Q-value $Q^*$ because an optimal policy can be then extracted by acting greedily with respect to it , i.e., $\pi^*(s)=\argmax_{a\in\mathcal{A}} Q^*(s,a)$.
In Q-learning, the learner's internal state includes only a Q-value estimate $Q:\mathcal{S}\times\mathcal{A}\to\R$ and therefore it is a value-based (and also model-free) RL algorithm.
The goal of Q-learning is to update $Q$ such that it eventually approximates $Q^*$ well and recommends the outputs the policy as $\widehat{\pi}(s) = \argmax_{a\in\mathcal{A}}Q(s,a)$.
{\em Tabular} Q-learning represents $Q$ as a table/vector with $|\mathcal{S}\times\mathcal{A}|$ entries, one entry per $(s,a)$ pair, and, upon a transition $(s_t,a_{t},r_t,s_{t+1})$, updates the entry of $(s_t,a_t)$ as
\begin{align}
\label{eq:Q-learning}
Q(s_t,a_t) \gets Q(s_t,a_t) + \alpha_t\left(r_t+\gamma\max_{a\in\mathcal{A}}Q(s_{t+1},a)-Q(s_t,a_t) \right). 
\end{align}
Here, the index $t$ is the total number of transitions the learner has experienced in Protocol \ref{protocol:RL interaction}, which keeps increasing across episodes;
$\alpha_t\in(0,1)$ controls how aggressive the update is, which in general depends on $t$.
It is instructive to compare update \eqref{eq:Q-learning} with {\em Q-value iteration}, the planning-setting counterpart summarized in Protocol \ref{protocol:Q_VI}, which computes $Q^*$ when the MDP $M=(\mathcal{S},\mathcal{A},P,R,\gamma)$ is fully known.
\begin{protocol}[htb]
\caption{Q-value iteration (planning)} \label{protocol:Q_VI}
\begin{algorithmic}[1]
\State planner is given $M=(\mathcal{S},\mathcal{A},P,R,\gamma)$ and initializes $Q(s,a)$ arbitrarily for all $(s,a)\in\mathcal{S}\times\mathcal{A}$;
\Repeat {~over iterations:}
    \For{every $(s,a)\in\mathcal{S}\times\mathcal{A}$} \Comment{{\color{Gray}a full sweep over all state-action pairs}}
        \State planner computes $Q'(s,a) \gets R(s,a)+\gamma\E_{s'\sim P(s,a)}\left[\max_{a'\in\mathcal{A}}Q(s',a')\right]$;
        \Comment{{\color{Gray}$(\mathcal{T}Q)(s,a)$, cf.~\eqref{eq:Bellman_optimality_operator}; requires known $(P,R)$}}
    \EndFor
    \State planner sets $Q \gets Q'$;
\Until{some stopping criterion is met}
\State planner outputs a policy $\widehat{\pi}$ as the greedy policy w.r.t.~$Q$;
\end{algorithmic}
\end{protocol}

Each iteration of Protocol \ref{protocol:Q_VI} overwrites {\em every} entry of $Q$ by an exact expectation under $(P,R)$, neither of which is available to the learner in the RL setting.
Tabular Q-learning can thus be viewed as a sample-based, incremental version of Protocol \ref{protocol:Q_VI}: it updates only the single entry $(s_t,a_t)$ that the learner just visited, replaces the expectation by the one-sample target \eqref{eq:Q-learning-target}, and, because a single sample is noisy, moves $Q(s_t,a_t)$ only partway toward that target with step size $\alpha_t$ instead of overwriting it.
Protocol \ref{protocol:Q_Learning} summarizes tabular Q-learning as an instance of Protocol \ref{protocol:RL interaction}.
\begin{protocol}[htb]
\caption{Tabular Q-learning algorithm} \label{protocol:Q_Learning}
\begin{algorithmic}[1]
\State learner initializes $Q(s,a)$ for all $(s,a)\in\mathcal{S}\times\mathcal{A}$ (e.g., $Q(s,a)=0$) and sets $t=1$;
\Repeat {~over episodes:}
    \State learner observes initial state $s\sim d_0$ of the current episode; 
    \For{timesteps within the episode}
        \State learner chooses action $a$ that is $\epsilon$-greedy w.r.t. $Q(s,\cdot)$;
        \Comment{{\color{Gray}exploration; $\epsilon$ usually decreases over time}}
        \State learner takes $a$ and  observes next state $s'$ and reward $r$;
        \State learner updates $Q(s,a) \gets Q(s,a) + \alpha_t\left(r+\gamma\max_{a'\in\mathcal{A}}Q(s',a')-Q(s,a) \right)$;
        \Comment{{\color{Gray}update \eqref{eq:Q-learning}}}
        \State learner sets $s\gets s'$ and $t\gets t+1$;
    \EndFor
\Until{some stopping criterion is met}
\State learner outputs a policy $\widehat{\pi}$ as the greedy policy w.r.t.~$Q$;
\end{algorithmic}
\end{protocol}

To understand why update \eqref{eq:Q-learning} makes sense, recall Q-value iteration: starting with an arbitrary $Q\in\R^{|\mathcal{S}\times\mathcal{A}|}$, update it iteratively as $Q \gets \mathcal{T} Q$ where {\em Bellman optimality operator} $\mathcal{T}:\R^{|\mathcal{S}\times\mathcal{A}|}\to\R^{|\mathcal{S}\times\mathcal{A}|}$ is defined as:
\begin{align}
  (\mathcal{T}Q)(s,a) :=&~ R(s,a)+\gamma\E_{s'\sim P(s,a)}\left[\max_{a\in\mathcal{A}}Q(s',a)\right] \label{eq:Bellman_optimality_operator}\\ 
  \approx &~r+\gamma\max_{a\in\mathcal{A}}Q(s',a) \label{eq:Q-learning-target} 
\end{align}
where the $\approx$ approximates the expectation by a single sample of transition $(s,a,r,s')$ with $r=R(s,a)$ and $s'\sim P(s,a)$.
In the planning setting, we can perform the Q-value iteration and we have the convergence of $\mathcal{T}^kQ\to Q^*$ as $k\in\infty$, where $\mathcal{T}^k$ iteratively applies $\mathcal{T}$ for $k$ times.
In the RL setting, we cannot perform a full value iteration update $(\mathcal{T}Q)$ because $(P,R)$ is unknown. Instead, upon on a new transition $(s,a,r,s')$, we update $Q(s,a)$ to be closer to the one-sample target \eqref{eq:Q-learning-target}:
\begin{align*}
    Q(s,a) \gets (1-\alpha) Q(s,a) + \alpha \left(r+\gamma\max_{a\in\mathcal{A}}Q(s',a)\right)
    = Q(s,a) + \alpha\left(r+\gamma\max_{a\in\mathcal{A}}Q(s',a)-Q(s,a)\right)
\end{align*}
where $\alpha\in(0,1)$ controls how aggressively $Q(s,a)$ is updated to target \eqref{eq:Q-learning-target}.
Setting $(s,a,r,s')=(s_t,a_t,r_t,s_{t+1})$ and $\alpha=\alpha_t$ in the update above recovers \eqref{eq:Q-learning}.

It turns out that, if the following conditions hold, we have the convergence guarantee of $Q \to Q^*$ as $t\to\infty$:  \begin{itemize}
    \item Every $(s,a)$ pair is visited infinitely often.
    \item $\alpha_t$ decreases but not too quickly.
\end{itemize}
A formal description of the first condition is 
$\lim_{t\to\infty}N_t(s,a)=\infty$ for any $(s,a)$, where $N_t(s,a)$ is the number of times $(s,a)$ is visited by $t$.
This means the learner must do {\em exploration} when choosing actions at line \ref{protocol:RL interaction:line:choose action} in Protocol \ref{protocol:RL interaction}. For example, one can do $\epsilon$-greedy w.r.t.~the current $Q$. 
A formal description of the second condition is $\sum_t \alpha_t=\infty, \sum_t \alpha_t^2<\infty$ (e.g., $\alpha_t=\frac{1}{t}$ satisfies this).

\section{Deep Q-Network}
The obvious issue of tabular Q-learning is the difficulty of scaling to large state/action spaces.
{\em Function approximation} solves this issue by representing the Q-value estimates with a generic function.
The intuition is that, by updating the Q-value on a specific state-action pair, the changes made to the function will also influence other state-action pairs.
\begin{protocol}[htb]
\caption{DQN algorithm} \label{protocol:DQN}
\begin{algorithmic}[1]
\State learner initializes parameter $\theta$ of the Q network and sets replay buffer $\mathcal{D}=\{\}$ and target network parameter $\bar{\theta} = \theta$;
\Repeat {~over episodes:}
    \State learner observes initial state $s\sim d_0$ of the current episode; 
    \For{timesteps within the episode}
        \State learner chooses action $a$ that is $\epsilon$-greedy w.r.t. $Q_{\theta}(s,\cdot)$;
        \Comment{{\color{Gray}$\epsilon$ usually decreases over time}}
        \State learner takes $a$ and  observes next state $s'$ and reward $r$;
        \State learner add $(s,a,r,s')$ to replay buffer $\mathcal{D}$; 
        \State learner samples a batch of $N$ transitions from $\mathcal{D}$;
        \State learner takes a gradient step minimizing $L(\theta)$;
        \State learner sets $\bar{\theta} \gets \theta$ every $c$ updates of $\theta$; \Comment{{\color{Gray}updates target network periodically}}
    \EndFor
\Until{some stopping criterion is met}
\State learner outputs a policy $\widehat{\pi}$ as the greedy policy w.r.t.~$Q_\theta$;
\end{algorithmic}
\end{protocol}

A representative work is {\em Deep Q-Network} (DQN), where the Q-value estimates is represented using a neural network, $Q_\theta:\mathcal{S}\times\mathcal{A}\to\R$ where $\theta$ is the neural network's parameters.
Although people tried neural networks for Q-learning before DQN, those prior efforts updated the parameters on the most recent transition, in a similar fashion as in tabular Q-learning, and observed limited success due to 
DQN improved the training (i.e., parameters updating) of the networks with the following innovations:
\begin{itemize}
    \item It stores the past transitions in a memory, often referred to as {\em replay buffer}, $\mathcal{D}=\{(s,a,r,s')\}$.
    \item On a given transition $(s,a,r,s')$, it forms target \eqref{eq:Q-learning-target} using another network $Q_{\bar{\theta}}$, which is of the same structure as $Q_{\theta}$ but with its parameter $\bar{\theta}$ updated more slowly than $\theta$.
    \item To update $\theta$, it samples a batch of transitions from the replay buffer and update $\theta$ with a single gradient step to reduce the Q-value estimates with the formed target values.
\end{itemize}
This leads to the following loss function for network $Q_\theta$:
\begin{align*}
    L(\theta) = \frac{1}{N}\sum_{i=1}^{N} \left( Q_\theta(s_i,a_i) - \big(r_i+\gamma \max_{a\in\mathcal{A}} Q_{\bar{\theta}}(s'_i, a)\big) \right)^2
    \quad \text{where }
    (s_i,a_i,r_i,s'_i)\sim\mathcal{D} \text{ for } i=1,\ldots,N.
\end{align*}